\documentclass[11pt]{article}

\usepackage[a4paper,margin=1in]{geometry}
\usepackage{amsmath,amssymb,amsthm}
\usepackage{enumitem}
\usepackage{xcolor}
\usepackage{fancyhdr}
\usepackage{needspace}
\usepackage{physics}

\definecolor{courseblue}{RGB}{0,0,255}

\setlength{\parindent}{0pt}
\setlength{\parskip}{0.35em}
\setlength{\headheight}{14pt}

\pagestyle{fancy}
\fancyhf{}
\fancyfoot[C]{\thepage}
\renewcommand{\headrulewidth}{0pt}

\renewcommand{\ip}[2]{\left\langle #1, #2 \right\rangle}

\newcounter{problem}
\newenvironment{problem}{%
    \refstepcounter{problem}%
    \Needspace{7\baselineskip}%
    \par\vspace{0.7em}%
    \noindent\textcolor{courseblue}{\rule{\textwidth}{0.4pt}}%
    \par\vspace{0.5em}%
    \begin{list}{}{%
        \setlength{\leftmargin}{2.3em}%
        \setlength{\labelwidth}{1.8em}%
        \setlength{\labelsep}{0.5em}%
        \setlength{\topsep}{0pt}%
        \setlength{\partopsep}{0pt}%
    }%
    \item[{[\theproblem]}]
}{%
    \end{list}
    \par\vspace{0.5em}
}

\newenvironment{parts}{%
    \begin{enumerate}[label={[\alph*]},leftmargin=2.2em,itemsep=0.65em,topsep=0.65em]
}{%
    \end{enumerate}
}

% Type each solution between \begin{answer} and \end{answer}.
\newenvironment{answer}{%
    \Needspace{4\baselineskip}%
    \par\vspace{0.35em}
    \noindent\textbf{Answer:}\par
}{%
    \par\vspace{0.25em}
}

\begin{document}

\begin{center}
    {\color{courseblue}\Large\textbf{AMATH/PMATH 331: Applied Real Analysis}}
    \hfill
    {\color{courseblue}\Large\textbf{Fall 2026}}

    \vspace{1.2em}
    {\color{courseblue}\Large
    Assignment 2 Submission; due Thursday, October 8, 2026}

    \vspace{1.2em}
    
    \textbf{Name:} Vlad Paun

    \textbf{Student number:} 21291453
\end{center}

\begin{problem}
Let $(\mathbf{x}_k)_{k\in\mathbb{N}}$ be a sequence in $\mathbb{R}^n$.
Write $\mathbf{x}_k=(x_{k1},\ldots,x_{kn})$, so that $(x_{kj})_{k\in\mathbb{N}}$
is a sequence in $\mathbb{R}$ for each $j=1,\ldots,n$.

\begin{parts}
    \Needspace{8\baselineskip}\item Let $\mathbf{a}=(a_1,\ldots,a_n)\in\mathbb{R}^n$.
    Prove that $\lim_{k\to\infty}\mathbf{x}_k=\mathbf{a}$ if and only if
    $\lim_{k\to\infty}x_{kj}=a_j$ for each $j=1,\ldots,n$.
    (\emph{Don't forget to prove both directions!})

    \begin{answer}
    \end{answer}

    \Needspace{8\baselineskip}\item Prove that $(\mathbf{x}_k)_{k\in\mathbb{N}}$
    is a Cauchy sequence in $\mathbb{R}^n$ if and only if
    $(x_{kj})_{k\in\mathbb{N}}$ is a Cauchy sequence in $\mathbb{R}$ for each
    $j=1,\ldots,n$. (\emph{Don't forget to prove both directions!})

    \begin{answer}
    \end{answer}
\end{parts}
\end{problem}

\begin{problem}
The two parts of this problem \emph{are} related.

\begin{parts}
    \Needspace{8\baselineskip}\item Let $E\subseteq\mathbb{R}$ be nonempty.
    If $E$ is bounded below, so that $\inf E$ exists, prove that there exists
    a sequence $(a_k)$ with $a_k\in E$ for all $k$, and
    $\lim_{k\to\infty}a_k=\inf E$.

    \begin{answer}
    \end{answer}

    \Needspace{12\baselineskip}\item Let $A\subseteq\mathbb{R}^n$ be nonempty,
    and $\mathbf{x}\in\mathbb{R}^n$. We define the \emph{distance} from $A$ to
    $\mathbf{x}$ to be
    \[
        \operatorname{dist}_A(\mathbf{x})
        =\inf\{\lVert\mathbf{x}-\mathbf{y}\rVert:\mathbf{y}\in A\}.
    \]
    (Note that $\operatorname{dist}_A(\mathbf{x})$ is well-defined and
    $\operatorname{dist}_A(\mathbf{x})\geq 0$, because it is the infimum of a
    nonempty set of nonnegative real numbers.) Prove that there exists a
    sequence $(\mathbf{y}_k)$ in $\mathbb{R}^n$ with $\mathbf{y}_k\in A$ for all
    $k$ and $\lim_{k\to\infty}\lVert\mathbf{x}-\mathbf{y}_k\rVert
    =\operatorname{dist}_A(\mathbf{x})$.

    \begin{answer}
    \end{answer}
\end{parts}
\end{problem}

\begin{problem}
Let $A\subseteq\mathbb{R}^n$. We define the \emph{closure} of $A$, denoted by
$\overline{A}$, to be the subset
\[
    \overline{A}=\{\mathbf{x}\in\mathbb{R}^n:
    B_\varepsilon(\mathbf{x})\cap A\neq\varnothing
    \text{ for all }\varepsilon>0\}.
\]
That is, a point $\mathbf{x}\in\mathbb{R}^n$ lies in the closure of $A$ if
and only if every open ball of any positive radius intersects $A$.

\begin{parts}
    \Needspace{6\baselineskip}\item Show that $A\subseteq\overline{A}$.

    \begin{answer}
    \end{answer}

    \Needspace{7\baselineskip}\item Show that
    $\overline{A}=\{\mathbf{x}\in\mathbb{R}^n:U\cap A\neq\varnothing
    \text{ for any open set }U\text{ containing }\mathbf{x}\}$.

    \begin{answer}
    \end{answer}

    \Needspace{9\baselineskip}\item Show that $\overline{A}$ is closed.

    \emph{Hint:} First show there exists $\varepsilon_0$ such that
    $B_{\varepsilon_0}(\mathbf{x})\cap A=\varnothing$, and then show that
    $B_{\varepsilon_0}(\mathbf{x})\cap\overline{A}=\varnothing$.
    You will need to use part [b].

    \begin{answer}
    \end{answer}

    \Needspace{6\baselineskip}\item If $A$ is closed, show that $\overline{A}=A$.

    \begin{answer}
    \end{answer}
\end{parts}
(Combining parts [c] and [d] shows that $A$ is closed if and only if
$\overline{A}=A$.)
\end{problem}

\begin{problem}
Let $E\subseteq\mathbb{R}^n$ be nonempty.

\begin{parts}
    \Needspace{8\baselineskip}\item If $(\mathbf{x}_k)_{k\in\mathbb{N}}$ is a
    sequence such that $\mathbf{x}_k\in E$ for all $k\in\mathbb{N}$, and
    $\lim_{k\to\infty}\mathbf{x}_k=\mathbf{p}$, prove that
    $\mathbf{p}\in\overline{E}$.

    \begin{answer}
    \end{answer}

    \Needspace{8\baselineskip}\item Let $\mathbf{p}\in\overline{E}$.
    Prove that there exists a sequence $(\mathbf{x}_k)_{k\in\mathbb{N}}$ such
    that $\mathbf{x}_k\in E$ for all $k\in\mathbb{N}$, and
    $\lim_{k\to\infty}\mathbf{x}_k=\mathbf{p}$.
    \emph{Hint:} Consider $B_{\frac{1}{k}}(\mathbf{p})$.

    \begin{answer}
    \end{answer}
\end{parts}
Taken together, these two statements say that $\overline{E}$ is exactly the
set of limits of convergent sequences whose elements all lie in $E$.
\end{problem}

\begin{problem}
Let $(\mathbf{x}_m)_{m\in\mathbb{N}}$ be a sequence in $\mathbb{R}^n$.
Define a new sequence $(\mathbf{s}_k)_{k\in\mathbb{N}}$ by
\[
    \mathbf{s}_k=\sum_{m=1}^k\mathbf{x}_m.
\]
If the sequence $(\mathbf{s}_k)$ converges, then we say that the \emph{series}
$\sum_{m=1}^{\infty}\mathbf{x}_m$ \emph{converges}.
If $\lim_{k\to\infty}\mathbf{s}_k=\mathbf{s}$, then we write
$\sum_{m=1}^{\infty}\mathbf{x}_m=\mathbf{s}$.

\begin{parts}
    \Needspace{12\baselineskip}\item Prove that
    $\sum_{m=1}^{\infty}\mathbf{x}_m$ converges iff: for all $\varepsilon>0$,
    there exists $N\in\mathbb{N}$ such that
    \[
        \left\lVert\sum_{m=k+1}^{\ell}\mathbf{x}_m\right\rVert<\varepsilon
    \]
    for all $\ell\geq k\geq N$.

    \emph{Hint:} A sequence in $\mathbb{R}^n$ converges if and only if it is Cauchy.

    \begin{answer}
    \end{answer}

    \Needspace{13\baselineskip}\item Suppose that
    $\sum_{m=1}^{\infty}\mathbf{x}_m$ converges. For any $N\in\mathbb{N}$,
    prove that $\sum_{m=N}^{\infty}\mathbf{x}_m$ converges, and that
    \[
        \lim_{N\to\infty}\left(\sum_{m=N}^{\infty}\mathbf{x}_m\right)=0.
    \]
    \emph{Hint:} Consider the sequence
    $(\tilde{\mathbf{s}}_\ell)_{\ell\geq N}$ given by
    $\tilde{\mathbf{s}}_\ell=\sum_{m=N}^{\ell}\mathbf{x}_m$.

    \begin{answer}
    \end{answer}
\end{parts}
\end{problem}

\begin{problem}
COMING SOON
\end{problem}

\begin{problem}
COMING SOON
\end{problem}

\begin{problem}
COMING SOON
\end{problem}

\begin{problem}
COMING SOON
\end{problem}

\begin{problem}
COMING SOON
\end{problem}

\end{document}
